\documentclass[11pt]{article}
\usepackage[margin=1in]{geometry}
\usepackage{amsmath, amssymb}

\title{Exercise 2.7: Unbiased Constant-Step-Size Trick}
\author{Sutton \& Barto, Section 2.5}
\date{}

\begin{document}
\maketitle

\section*{Problem}

In most of this chapter we have used sample averages to estimate action values because
sample averages do not produce the initial bias that constant step sizes do.
However, sample averages perform poorly on nonstationary problems.
Is it possible to avoid the bias of constant step sizes while retaining their advantages
on nonstationary problems? One way is to use a step size of
\begin{equation}
    \beta_n = \frac{\alpha}{\bar{o}_n} \tag{2.8}
\end{equation}
to process the $n$th reward for a particular action, where $\alpha > 0$ is a conventional
constant step size, and $\bar{o}_n$ is a trace of one that starts at 0:
\begin{equation}
    \bar{o}_n = \bar{o}_{n-1} + \alpha(1 - \bar{o}_{n-1}), \quad
    \text{for } n > 0, \quad \text{with } \bar{o}_0 = 0. \tag{2.9}
\end{equation}
Carry out an analysis like that in (2.6) to show that $Q_n$ is an exponential
recency-weighted average \emph{without initial bias}.

\section*{Solution}

\subsection*{Step 1: Solve the recurrence for $\bar{o}_n$}

Rewrite the recurrence as $\bar{o}_n = \alpha + (1-\alpha)\bar{o}_{n-1}$.
Unrolling with $\bar{o}_0 = 0$:
\begin{align*}
    \bar{o}_1 &= \alpha \\
    \bar{o}_2 &= \alpha + (1-\alpha)\alpha = \alpha\bigl[1 + (1-\alpha)\bigr] \\
    \bar{o}_3 &= \alpha\bigl[1 + (1-\alpha) + (1-\alpha)^2\bigr]
\end{align*}
In general, this is a geometric series:
\[
    \boxed{\bar{o}_n = \alpha \sum_{i=0}^{n-1}(1-\alpha)^i = 1 - (1-\alpha)^n}
\]
Note that $\bar{o}_n \to 1$ as $n \to \infty$, and $\bar{o}_n < 1$ for all finite $n$.

\subsection*{Step 2: Compute $\beta_n$}

\[
    \beta_n = \frac{\alpha}{\bar{o}_n} = \frac{\alpha}{1 - (1-\alpha)^n}
\]

\subsection*{Step 3: Expand $Q_{n+1}$ as a weighted sum}

The update rule $Q_{n+1} = Q_n + \beta_n(R_n - Q_n) = (1 - \beta_n)Q_n + \beta_n R_n$
expands recursively to:
\[
    Q_{n+1} = \underbrace{\prod_{j=1}^{n}(1-\beta_j)}_{w_0} \cdot Q_1
    \;+\; \sum_{i=1}^{n} \underbrace{\beta_i \prod_{j=i+1}^{n}(1-\beta_j)}_{w_i} \cdot R_i
\]

\subsection*{Step 4: Show $w_0 = 0$ (no initial bias)}

We need to evaluate $\prod_{j=1}^{n}(1-\beta_j)$. First, compute $1 - \beta_j$:
\[
    1 - \beta_j = 1 - \frac{\alpha}{\bar{o}_j} = \frac{\bar{o}_j - \alpha}{\bar{o}_j}
\]
From the recurrence, $\bar{o}_j = \alpha + (1-\alpha)\bar{o}_{j-1}$, so
$\bar{o}_j - \alpha = (1-\alpha)\bar{o}_{j-1}$. Therefore:
\[
    1 - \beta_j = \frac{(1-\alpha)\,\bar{o}_{j-1}}{\bar{o}_j}
\]
The product \textbf{telescopes}:
\[
    \prod_{j=1}^{n}(1-\beta_j)
    = \prod_{j=1}^{n} \frac{(1-\alpha)\,\bar{o}_{j-1}}{\bar{o}_j}
    = (1-\alpha)^n \cdot \frac{\bar{o}_0}{\bar{o}_n}
\]
Since $\bar{o}_0 = 0$:
\[
    \boxed{w_0 = \prod_{j=1}^{n}(1-\beta_j) = 0}
\]

Therefore $Q_{n+1}$ depends only on $R_1, \ldots, R_n$ --- the initial estimate $Q_1$
has \textbf{zero weight}. There is no initial bias.

\subsection*{Step 5: Verify the weights sum to 1}

Since $w_0 = 0$, the weights on $R_1, \ldots, R_n$ must sum to 1.
This follows from the same argument as for constant $\alpha$ in equation~(2.6):
the update rule preserves $\sum w_i = 1$, but now $w_0 = 0$ instead of $(1-\alpha)^n$.

\subsection*{Summary}

The trick works because:
\begin{itemize}
    \item $\bar{o}_n$ tracks a ``trace of one'' --- it starts at 0 and approaches 1,
          measuring how much data has been seen.
    \item $\beta_n = \alpha/\bar{o}_n$ is large early on (when $\bar{o}_n$ is small)
          and approaches $\alpha$ as $n \to \infty$.
    \item The telescoping product $\prod(1-\beta_j)$ contains $\bar{o}_0 = 0$ as a factor,
          killing the dependence on $Q_1$ entirely.
    \item In the long run, $\beta_n \approx \alpha$, so the method behaves like a constant
          step-size method and tracks nonstationary environments --- but without the
          initial bias.
\end{itemize}

\end{document}
